% Existing Short Summary from summary.tex, preserved verbatim.
% Author: review this draft against the final scope and findings; maximum 300 words.
Speech models such as text-to-speech and automatic speech recognition typically operate on discretised representations of both acoustic time and computational depth. Although speech is a continuous physical signal, it is usually represented as fixed-rate frames or tokens and processed by a fixed sequence of neural-network layers. In conditional generative speech models, text is also commonly introduced through auxiliary conditioning channels while the source distribution remains an abstract Gaussian, making the interpretation of the resulting conditional likelihood unclear. Continuous-depth models have nevertheless shown that neural transformations can be represented as invertible dynamical systems and evaluated using the change-of-variables formula.

Here, we investigate speech modelling in which both acoustic time and computational depth are continuous. This formulation allows signals to be evaluated at arbitrary temporal resolutions, enabling variable-length utterances to be packed without padding and permitting resolution to increase progressively over depth from sparse linguistic inputs to dense acoustic outputs. More importantly, it creates a route towards defining an invertible transformation from a continuous representation of text to a continuous representation of speech, so that likelihood describes transport between the conditioning and target distributions more directly. Together, continuous time and continuous depth provide a unified framework for adaptive temporal resolution, coarse-to-fine alignment, efficient computation, and probabilistically grounded mappings between language and speech.
